\subsection{An actual metabelian quotient and its coherent norm problem}

\begin{statement}{A.77}{Proved}
{Every finite root-torsor norm test for one fixed ordinary pair}
Fix an odd prime $\ell$.  Let $(A,B)$ be an ordinary completed
unit-equation pair over $\Q$ with $v_3(A)=1$ in $\widehat\Z$.
Write $a,b$ for continuous cocycles representing its coefficient-$\ell$
Kummer classes, in the tangential framing of A.75.
There is an actual quotient of the full geometric pro-$\ell$
metabelian group,
\[
 P=F/[N,N],\qquad N=\ker(F\longrightarrow\Z_\ell(1)),
 \qquad x\longmapsto1,\quad y\longmapsto0,
 \tag{A.77.1}
\]
with the commutator subgroup closed.  For every $n\ge1$,
the pair $(a,b)$ lifts to the arithmetic extension with geometric group
\[
 P_n=\Z_\ell[\Z/\ell^n\Z]\rtimes\Z_\ell,
 \tag{A.77.2}
\]
obtained from $P$ by the quotient map specified below.
This is one fixed ordinary pair for all these quotients.
The same conclusion, with the same norm-tower criterion below,
holds at $\ell=2$ for every such pair with $A/3$ a square
in $\widehat{\Q^\times}$.  In particular the constructed
point-$3$ congruent ghosts pass all these individual norm
levels at every coefficient prime.

An actual lift to $P$ itself is equivalent to a compatible sequence
of completed norm preimages in an explicit nested tower of number
fields $K_n$.  The exact compatibility class belongs to
\[
 \varprojlim{}^1_n
 \ker\left(N_{K_n/\Q}:
     \widehat{K_n^\times}_{\ell}\longrightarrow
     \widehat{\Q^\times}_{\ell}\right).
 \tag{A.77.3}
\]
Here the completions are the inverse limits for $\ell$-power
exponents.  We compute the class and its equivalence with lifting;
we do not assert that it vanishes.  The map from $F/F''$ to $P$
and its kernel are included in the proof.
\end{statement}
\begin{proof}
\emph{The geometric quotient and its original arithmetic action.}
The pro-$\ell$ Schreier generators of $N$ are the conjugates of
$y$ by the powers of $x$.  Their abelianization gives
\[
 N/[N,N]=\Lambda_X y,\qquad
 \Lambda_X=\Z_\ell[[X-1]],\qquad
 P=(\Lambda_X y)\rtimes\Z_\ell x.
 \tag{A.77.4}
\]
This identification also follows directly from the Fox map in A.75:
specialize $Y$ to one.  The commutator $[x,y]$ maps to
$(X-1)y$, and the full derived module
$\Z_\ell[[X-1,Y-1]]e$ maps by
$h(X,Y)e\mapsto h(X,1)(X-1)y$.
Its kernel is exactly $(Y-1)D$, since $X-1$ is a nonzero divisor.
The resulting group is generated by $x,y$ with the displayed
abelian normal kernel; the same continuous Fox map proves that
there is no further kernel.  Thus the quotient map is explicitly
\[
 F/F''\longrightarrow P,
 \qquad \ker=(Y-1)D.
 \tag{A.77.5}
\]
All maps are continuous homomorphisms of compact pro-$\ell$ groups.

The tangential action has
$g(x)=x^{\chi_\ell(g)}$ and
$g(y)=f_g^{-1}y^{\chi_\ell(g)}f_g$, where $f_g\in F'\subset N$.
In $P$ the images of $f_g$ and $y$ commute.  Therefore its actual
action is
$g(x)=x^{\chi_\ell(g)}$, $g(y)=y^{\chi_\ell(g)}$.
It sends $h(X)y$ to
$\chi_\ell(g)h(X^{\chi_\ell(g)})y$.
Twisting by $a$ gives the compact coefficient module
\[
 \Lambda_A(1)=\Lambda_X,
 \qquad g\cdot h(X)=
  \chi_\ell(g)X^{a(g)}h(X^{\chi_\ell(g)}).
 \tag{A.77.6}
\]
The group law in (A.77.4) now shows directly that a lift of
$(a,b)$ is a continuous one-cocycle in $\Lambda_A(1)$ with
augmentation $b$.  Equality of augmentation classes suffices:
an augmentation coboundary can be removed by lifting its
zero-cochain along the surjective augmentation of modules.
The finite quotient (A.77.2) sets $X^{\ell^n}=1$ in the
normal module and retains the full $\Z_\ell$ coordinate of $x$.
The ideals are stable under (A.77.6), so these are actual
arithmetic quotient maps over $G_\Q$.

\emph{The fields and the precise norm comparison.}
Choose positive rational representatives $a_n$ of $A$ modulo
$\ell^n$ with $v_3(a_n)=1$ exactly.  Such representatives exist:
adjust the valuation by a rational $\ell^n$th power, and the
sign by $(-1)^{\ell^n}$ when necessary.
Compatibility gives unique positive rational $d_n$ with
$a_{n+1}/a_n=d_n^{\ell^n}$.
Choose the positive real roots and put
\[
 \alpha_n^{\ell^n}=a_n,\qquad
 K_n=\Q(\alpha_n),\qquad
 \alpha_{n+1}^{\ell}=d_n\alpha_n.
 \tag{A.77.7}
\]
The polynomials are Eisenstein over $\Q_3$, so
$[K_n:\Q]=\ell^n$; (A.77.7) gives the stated nested field maps.
This includes $\ell=3$.

The $G_\Q$-set of roots of $Z^{\ell^n}-a_n$ has the affine
action $j\mapsto a(g)+\chi_\ell(g)j$ after a compatible choice
of roots of unity and Kummer cocycles.  Therefore
\[
 \Lambda_A(1)/(X^{\ell^n}-1)
   =\Z_\ell[\operatorname{Roots}(Z^{\ell^n}-a_n)](1).
 \tag{A.77.8}
\]
Changing cocycle representatives conjugates this identification
by a translation of the root torsors.
Continuous Shapiro and Kummer theory identify its $H^1$ with
$\widehat{K_n^\times}_{\ell}$.
The augmentation map is corestriction to $G_\Q$, hence the field
norm; the transition sending each root to its image in (A.77.7)
is corestriction from $G_{K_{n+1}}$ to $G_{K_n}$, hence
$N_{K_{n+1}/K_n}$.
These identifications follow already on the finite Kummer
sequences by multiplying the conjugate roots, and commute with
their power transitions.

The modules in (A.77.8) are finite free over $\Z_\ell$;
their transition maps have continuous linear sections obtained
by choosing one lift of each root.  Continuous cochains are
therefore degreewise surjective under transition.
Furthermore
\[
 H^0(G_\Q,\Z_\ell[\operatorname{Roots}(Z^{\ell^n}-a_n)](1))
       =H^0(G_{K_n},\Z_\ell(1))=0:
 \tag{A.77.9}
\]
a number field contains only finitely many roots of unity, and
has no nonzero compatible system of roots of all $\ell$-power
orders.  The continuous cochain inverse-limit sequence gives
the exact isomorphism
\[
 H^1(G_\Q,\Lambda_A(1))
   \xrightarrow{\ \sim\ }
 \varprojlim_{n,N_{K_{n+1}/K_n}}
             \widehat{K_n^\times}_{\ell}.
 \tag{A.77.10}
\]
This proves the asserted coherent-norm interpretation of $P$.

\emph{Every individual norm test is satisfied.}
Let $m=\ell^n$ and $L_n=\Q(\mu_m,\alpha_n)$.
We verify the full group needed for the norm theorem.
The number $a_n$ is not an $\ell$th power in $\Q(\mu_m)$.
Indeed such a power would place its real $\ell$th root in this
abelian extension, since all $\ell$th roots of unity are present.
The field of that real root would be an abelian Galois extension
of $\Q$, and therefore all its conjugates would be real.
For odd $\ell$, only one root of $Z^\ell-a_n$ is real;
its degree would be one, contradicting $v_3(a_n)=1$.
Kummer theory now gives
$[L_n:\Q(\mu_m)]=m$.
The faithful action on the roots and on $\mu_m$ embeds its
Galois group in the affine group with exactly the same order.
Consequently
\[
 \operatorname{Gal}(L_n/\Q)
   = C_m\rtimes(\Z/m\Z)^\times,
 \qquad (\Z/m\Z)^\times=C_{\ell^{n-1}(\ell-1)}.
 \tag{A.77.11}
\]
This group is metacyclic and its Schur multiplier is zero.
Here is the calculation of the latter assertion.  In the
homology spectral sequence for its cyclic normal subgroup,
the terms of total degree two are
$H_2(C_{\ell^{n-1}(\ell-1)},\Z)=0$,
$H_1(C_{\ell^{n-1}(\ell-1)},C_m)=0$, and
$H_0(C_{\ell^{n-1}(\ell-1)},H_2(C_m,\Z))=0$.
The middle equality follows from the cyclic resolution:
the generator acts by a primitive unit $r$ modulo $m$,
and $r-1$ is invertible modulo $\ell$.
Thus $H_2$ of the group is zero, which is its Schur multiplier.

Hoshi--Yamasaki, Theorems 1.2 and 3.2, applies to a finite
separable extension of a global field whose Galois closure
has precisely this metacyclic group and trivial Schur multiplier;
it proves the Hasse norm principle for $K_n/\Q$
\cite{RootHoshiYamasaki2025}.
Choose a rational representative $b_n$ of $B$ modulo $m$.
At each finite place $p$ the ordinary pair supplies
$x_p\ne0,1$ with $A_p=\kappa(x_p)$ and
$B_p=\kappa(1-x_p)$.
There is $z_p\in\Q_p^\times$ with $x_p=a_nz_p^m$, and
\[
 N_{K_n\otimes\Q_p/\Q_p}(1-z_p\alpha_n)=1-x_p.
 \tag{A.77.12}
\]
The ratio $b_n/(1-x_p)$ is an $m$th power in $\Q_p^\times$,
and is the norm of that scalar root in the degree-$m$ etale
algebra.  Thus $b_n$ is an exact local norm at every finite
place.  At the real place every real scalar is an $m$th
power because $m$ is odd, so it is a norm there as well.
At almost all finite places the extension is unramified and
$b_n$ is a unit; the norm on the units of an unramified local
field extension is surjective.  These local preimages can
therefore be taken to be units at almost all places, giving
an idele norm as required by the theorem.
We obtain $c_n\in K_n^\times$ with $N_{K_n/\Q}(c_n)=b_n$.

The kernel of reduction modulo $m$ on
$\widehat{\Q^\times}_{\ell}$ is $m$ times that completion.
This follows directly by dividing the prime valuations by
$m$ in the cofinally convergent prime lattice; $\ell$ is odd,
so the sign contributes no factor.
Choose $D_n$ with $D_n^m=B/b_n$ and regard it as a scalar
in $\widehat{K_n^\times}_{\ell}$.
Then $\beta_n=c_nD_n$ satisfies $N_{K_n/\Q}(\beta_n)=B$.
Together with (A.77.8) this constructs the stated lift to
every $P_n$ for the same pair.

\emph{The exact remaining compatibility class.}
Put $Z_n=\ker N_{K_n/\Q}$ on these completions.
For any choices $\beta_n$ just constructed, set
\[
 d_n'=\beta_n\,N_{K_{n+1}/K_n}(\beta_{n+1})^{-1}\in Z_n.
 \tag{A.77.13}
\]
In the countable inverse system of abelian groups, define
$\varprojlim^1 Z_n$ as the cokernel of the map on
$\prod_n Z_n$ sending $(z_n)$ to
$(z_nN_{K_{n+1}/K_n}(z_{n+1})^{-1})$.
Changing the choices $\beta_n$ changes (A.77.13) by the image
of exactly this map.  Its class is therefore well-defined.
It is zero exactly when multiplying the $\beta_n$ by elements
of $Z_n$ makes them norm compatible.  Equations
(A.77.6)--(A.77.10) then prove, in both directions, equivalence
with an actual continuous lift to $P$.
The fields' completed multiplicative groups contain cofinally
convergent prime lattices with infinite discrete finite-exponent
quotients.  We have used their actual cochain inverse limit;
we have supplied no compactness argument that would set this
class to zero.
\end{proof}

% Unnumbered insertion in A.77 before its final certificate paragraph.
\paragraph{The dyadic norm levels for the constructed pairs.}
The same fixed-level norm conclusion holds at $\ell=2$
for every ordinary pair under consideration with
$v_3(A)=1$ and $A/3$ a square in
$\widehat{\Q^\times}$.
In particular it applies to the point-$3$ congruent
ghosts constructed above.
Use the compatible rational representatives $a_n$
of A.77, with $m=2^n$.
Each $a_n$ has the same rational square class as $3$,
so it has the same $\Q_2$ square class as $3$.

The three quadratic subfields of
$\Q_2(\mu_{2^\infty})$ are
$\Q_2(i)$, $\Q_2(\sqrt2)$ and
$\Q_2(\sqrt{-2})$.
Indeed its Galois group is
$\Z_2^\times\simeq C_2\times\Z_2$; its three
index-two subgroups already occur in
$\Q_2(\mu_8)$.
The square class of $3$ differs from
$-1,2,-2$ in $\Q_2^\times/\Q_2^{\times2}$:
the last two have odd valuation, and
$3/(-1)\equiv5\pmod8$.
It follows that $a_n$ is not a square in
$\Q_2(\mu_m)$.
Kummer theory over this field therefore gives
\[
 [\Q_2(\mu_m,a_n^{1/m}):\Q_2(\mu_m)]=m .
\]
The cyclotomic extension at two itself has degree
$\varphi(m)$.
For the global splitting field
$L_n=\Q(\mu_m,a_n^{1/m})$ this proves
that the decomposition group at two has order
$m\varphi(m)$, the largest possible global order.
Consequently it is the whole group
$G_n=\operatorname{Gal}(L_n/\Q)$.

Here is the exact Hasse-principle implication, including
the inflation map.
For any number-field torus $T$ split by a finite Galois
extension $L/k$, let $X=X^*(T)$ and $G=\operatorname{Gal}(L/k)$.
Tate--Poitou duality identifies the norm-torus
Tate--Shafarevich group with the dual of the kernel of
\[
 H^2(G,X)\longrightarrow
                  \prod_v H^2(G_v,X).
\]
This is the general statement in
\cite[Proposition 6.4 and Corollary 6.5, p.~35]
{BIR_Oki2026Norm}; it is not restricted to the
prime-squared extensions studied elsewhere in that paper.
To verify the finite-group comparison directly, a class
in $\operatorname{Sha}^2(k,X)$ restricts over $L$ to a tuple of
characters, because
$H^2(G_L,\Z^r)=\operatorname{Hom}_{\rm cts}(G_L,\Q/\Z)^r$.
All these characters are locally trivial, hence vanish
by Chebotarev.
Moreover $H^1(G_L,X)=0$, since a continuous map from
a compact group to the discrete torsion-free lattice
$X$ has finite, therefore zero, image.
The Hochschild--Serre sequence thus puts the class
in the image of inflation from $H^2(G,X)$.
At a place with $G_v=G$, the same vanishing
$H^1(G_{L_w},X)=0$ makes local inflation injective.
Local triviality at that one place therefore kills
the inflated global class.
Duality proves $\operatorname{Sha}^1(k,T)=0$.
For $T=R^1_{K/k}\G_m$, Hilbert 90 identifies
$H^1(k,T)=k^\times/N_{K/k}K^\times$ and the local
maps with their ordinary norm maps.
This proves the Hasse norm principle for $K/k$.

Apply this argument to $K_n/\Q$ and its splitting
field $L_n$, with the full decomposition group at two
just proved.
The finite-place local norm identity (A.77.12) is
unchanged for $m=2^n$.
The rational number $a_n$ is positive, so $K_n$ has
a real embedding; the norm of its real etale algebra
onto $\R^\times$ is surjective.
Thus a rational representative $b_n$ of $B$ modulo
$m$ is a local norm at every place and is a global
norm.
Finally, the kernel of reduction modulo $2^n$ on
$\widehat{\Q^\times}_2$ is exactly its subgroup of
$2^n$th powers: the sign is already detected modulo
$2^n$, and the remaining prime coordinates can be
divided by $2^n$ with the same cofinal convergence.
The scalar correction $D_n$ in A.77 therefore gives
the actual completed norm preimage of $B$ also at two.
The remaining norm-tower compatibility class is
exactly the one already displayed in (A.77.13).


The finite certificate \path{code/norm_tower_certificate.py}
checks the norm polynomial identities, the finite affine groups,
the cyclic homology vanishing inputs, and the coefficient
transition maps.  The proof supplies the global norm theorem's
hypotheses and the continuous comparisons.  A.77 does not give
a lift to $F/F''$: the exact comparison kernel remains
$(Y-1)D$ in (A.77.5).
